     1	\ifx\danslelivre\undefined
     2	\RequirePackage{luatex85}
     3	\documentclass[10pt]{../fichiers-configuration-latex/smfart}
     4	\input{../fichiers-configuration-latex/paquets}
     5	\input{../fichiers-configuration-latex/numerotation}
     6	\input{../fichiers-configuration-latex/macros}
     7	
     8	
     9	\externaldocument{../01/01.tex}
    10	\externaldocument{../02/02.tex}
    11	\externaldocument{../03/03.tex}
    12	\externaldocument{../04/04.tex}
    13	\externaldocument{../05/05.tex}
    14	\externaldocument{../05bis/05bis.tex}
    15	\externaldocument{../06/06.tex}
    16	\externaldocument{../07/07.tex}
    17	\externaldocument{../08/08.tex}
    18	\externaldocument{../09/09.tex}
    19	%\externaldocument{../10/10.tex}
    20	\externaldocument{../11/11.tex}
    21	\externaldocument{../12/12.tex}
    22	\externaldocument{../13/13.tex}
    23	\externaldocument{../14/14.tex}
    24	\externaldocument{../15/15.tex}
    25	\externaldocument{../16/16.tex}
    26	\externaldocument{../17/17.tex}
    27	\externaldocument{../18/18.tex}
    28	\externaldocument{../19/19.tex}
    29	
    30	\setboolean{original}{false}
    31	
    32	\begin{document}
    33	
    34	\begin{center}
    35	XI. Comparaison avec la cohomologie classique : cas d'un schema lisse\\
    36	M. Artin 
    37	\end{center}
    38	
    39	\selectlanguage{english}version : \showgitstatus\selectlanguage{french}
    40	\tableofcontents
    41	
    42	\else
    43	\chapter{Comparaison avec la cohomologie classique : cas d'un schema lisse}
    44	%\addtocontents{toc}{. }
    45	\begin{center}
    46	M. Artin 
    47	\end{center}
    48	\fi
    49	
    50	\setcounter{pageoriginal}{63}
    51	
    52	\section{Introduction}\etiquette[1]{XI.sec1}\pageoriginale
    53	
    54	Dans le nᵒ \hyperref[XI.sec4]{4} nous démontrerons que la cohomologie étale
    55	à valeurs dans un faisceau localement constant de torsion coïncide
    56	avec la cohomologie classique pour un schéma $X$ lisse sur
    57	$\Spec\CC$. La démonstration, qui est essentiellement celle
    58	donnée dans \cite{1}, est élémentaire à partir du théorème de Grauert
    59	rappelé plus bas (\hyperref[XI.thm4.3]{4.3} (iii)), qui dit que chaque revêtement fini de
    60	$X(\CC)$ provient d'un revêtement étale de $X$; c'est pourquoi
    61	on la donne ici, bien qu'on prouvera un résultat plus complet plus
    62	tard, en utilisant le théorème de résolution des singularités de
    63	Hironaka \cite{3} \footnote[1]{Notons d'ailleurs que la résolution des
    64	  singularités de Hironaka permet également de donner une
    65	  démonstration très simple du résultat de Grauert, et à ce titre peut
    66	être considérée comme étant de toutes façons à la base des
    67	théorèmes de comparaison. Voir à ce sujet l'exposé de Mme M. Raynaud
    68	dans \SGA 1 XII (North Holland Pub. Cie).}. De plus, les bons
    69	voisinages  construits dans la section 3 pourraient être utiles dans
    70	d'autres\pageoriginale situations.
    71	
    72	Notons qu'il n'est pas vrai que la cohomologie étale et la cohomologie
    73	classique coïncident pour les faisceaux qui ne sont pas de
    74	torsion. Par exemple, on a $H^1(X_{\et},\ZZ)=0$ pour tout
    75	schéma normal $X$ (\hyperrefext[IX][IX.prop3.6]{3.6} (ii)), mais la cohomologie classique n'est
    76	pas nécessairement nulle.
    77	
    78	\section{Existence de sections hyperplanes assez générales}\etiquette[2]{XI.sec2}
    79	
    80	Le résultat suivant est classique mais on manque de référence (en
    81	attendant \EGA V):
    82	
    83	\begin{enonce*}{Théorème 2.1}\etiquette[2.1]{XI.thm2.1} %Pourquoi l'énoncé
    84	                                %n'est pas formulé pour un sous-espace
    85	                                %linéaire qui tourne autour d'un point
    86	                                %fixé? c'est-à-dire (i) avec $P\in L$ ou
    87	                                %$(ii)$ avec $d=1$
    88	Soit $k$ un corps et $V\subset \PP^n_k$ un schéma algébrique
    89	projectif purement de dimension $r$. Soit $V'$ le sous-schéma ouvert
    90	des points de $V$ lisses sur $\Spec k$ et soit $P$ un point rationnel
    91	de $\PP^n$. Alors:
    92	\begin{enumerate}
    93	\renewcommand{\theenumi}{\roman{enumi}}
    94	\renewcommand{\labelenumi}{(\theenumi)}
    95	\item Un sous-espace linéaire « assez général » $L$ de dimension
    96	  donnée $d$ de $\PP^n$ coupe $V$ transversalement en chaque
    97	  point de $L\cap V'$.
    98	
    99	\item Soit $\{H_1,\ldots,H_s\}$ $(s\leq r)$ un système d'hyperplans de %un
   100	                                %'hyperplan' n'est pas linéaire ici
   101	  degré $d\geq 2$ passant par $P$ et soit $Y=H_1\cap\ldots \cap
   102	  H_s$. Si $\{H_1,\ldots,H_s\}$ est assez général, alors $Y$ coupe $V$
   103	  transversalement en chaque point de $Y\cap V'$.
   104	\end{enumerate}
   105	\end{enonce*}
   106	
   107	La locution « pour un sous-espace linéaire assez général » veut
   108	dire qu'il existe un sous-ensemble ouvert dense $U$ de la
   109	grassmannienne qui paramétrise les sous-espaces linéaires de dimension % 'qui paramètre'
   110	$d$, tel que l'assertion envisagée sur $L$ soit vraie chaque fois que
   111	le paramètre de $L$ est dans $U$. Cela n'implique donc pas qu'il
   112	existe un tel $L$ (défini sur $k$), si $k$ n'est pas défini. De même,
   113	% mentionner le plongement de Veronese?
   114	les hyperplans de degré $d$ passant par $P$ sont 
   115	paramétrisés\pageoriginale par un espace projectif $\PP^N$, et la %'paramétrés'
   116	locution « pour $\{H_1,\ldots,H_s\}$ assez général » veut dire
   117	qu'il existe un sous-ensemble ouvert dense $U$ de $(\PP^N)^s$
   118	tel que l'assertion envisagée soit vraie chaque fois que le paramètre
   119	de $(H_1,\ldots,H_s)$ est dans $U$.
   120	
   121	Démonstration (i): Par récurrence on est ramené immédiatement au cas
   122	où $L$ est un hyperplan. Soit $(V'_{\nu})$ un recouvrement fini ouvert de
   123	$V'$. Si l'assertion (i) est vraie quand on remplace le symbole $V'$
   124	par $V'_{\nu}$ quel que soit $\nu$, c'est aussi vrai pour $V'$, parce
   125	que l'intersection d'un nombre fini de sous-ensemble ouverts denses
   126	est également ouvert et dense. On peut donc remplacer la situation par
   127	un schéma algébrique affine, disons $V\subset \EE^n_k=\Spec
   128	k[x_1,\ldots,x_n]$. L'hyperplan $L$ sera l'ensemble des zéros d'une
   129	équation linéaire
   130	$$
   131	\ell(x)=a_0+\sum^n_{i=1}a_ix_i{\quoi}.
   132	$$
   133	
   134	Rappelons qu'un schéma lisse est localement intersection complète. On
   135	peut donc remplacer $V'$ par un sous-ensemble ouvert de points de $V$
   136	tel qu'il existe des polynômes $f_1(x),\ldots,f_{n-r}(x)$ qui
   137	s'annulent sur $V$ et des indices $1\leq j_{\nu}\leq n$
   138	$(\nu=1,\ldots,n-r)$ tels que le déterminant
   139	$$
   140	\left|\frac{\partial f_i}{\partial x_{j_{\nu}}}\right|\quad 
   141	(i,\nu=1,\ldots,n-r)
   142	$$
   143	soit inversible. Disons $j_{\nu}=\nu$ $(\nu=1,\ldots,n-r)$. Alors on peut
   144	résoudre les équations
   145	$$
   146	y_j=\sum^{n-r}_{i=1}\quad c_i\quad \frac{\partial f_i}{\partial x_j}
   147	$$
   148	pour\pageoriginale $c_i$ comme fonction rationnelle des variables $x_i$, $y_j$
   149	$(i=1,\ldots, n-r; j=1,\ldots,n)$, et ces fonctions $c_i(x,y)$ sont
   150	définies en chaque point de $V'$.
   151	
   152	Or dire que $L$ coupe $V$ transversalement en un point $Q$ de $V'$
   153	revient à dire que la valeur de la matrice
   154	$$
   155	\begin{bmatrix}
   156	\frac{\partial f_1}{\partial x_1} & \cdots & \frac{\partial
   157	  f_1}{\partial x_n}\\
   158	\vdots & &\vdots\\
   159	\frac{\partial f_{n-r}}{\partial x_1} & & \frac{\partial
   160	  f_{n-r}}{\partial x_n}\\
   161	a_1 & \cdots & a_n
   162	\end{bmatrix}
   163	$$
   164	en $Q$ est de rang $n-r+1$. Donc pour que $L$ coupe $V$
   165	transversalement en chaque point de $V'$ il faut et suffit que le
   166	système des $r+1$ équations
   167	\begin{align*}
   168	a_j &= \sum^{n-r}_{i=1} c_i (x,a)\frac{\partial f_i}{\partial x_j}\qquad
   169	(j=n-r+1,\ldots, n)\\
   170	a_0 &+ \sum^n_{i=1} a_i x_i =0
   171	\end{align*}
   172	n'ait aucune solution sur $V'$. Une telle condition est évidemment
   173	constructible en $(a)$: Soit $X=V'\times A$, où $A$ est l'espace affine
   174	de coordonnées $a_0,\ldots,a_n$, et soit $Y\subset X$ le
   175	sous-ensemble fermé des solutions de ces équations. La condition est
   176	que la fibre de $Y/A$ soit vide, ce qui est bien une
   177	condition\pageoriginale constructible (\EGA IV 9.5.1). Il suffit donc % \ref
   178	de vérifier que la fibre au-dessus du point générique de $A$ est vide,
   179	ce qui se voit immédiatement sur la forme des équations. Ceci achève
   180	la démonstration de (i). (ii): Considérons l'application rationnelle
   181	$\PP^n\rightarrow \PP^N$ donnée par le système linéaire
   182	des hyperplans de degré $d$ ($d\geq 2$) passant par $P$. Il est bien
   183	connu (et on le vérifie aisément) que c'est un « éclatement » de
   184	$P$, donc donne une immersion localement fermée
   185	$$
   186	f:\PP^n-\{P\}\longrightarrow \PP^N{\quoi}.
   187	$$
   188	Soit $V''=V'-(\{P\}\cap V')$, $\bar{V}{}''=\sisi{f}{F}(V'')$ et soit $L_i$
   189	l'hyperplan de $\PP^N$ correspondant aux hyperplans $H_i$
   190	$(i=1,\ldots,s)$. Alors $V''$ est isomorphe à $\bar{V}{}''$, et d'après
   191	la définition de $f$, dire que $Y$ coupe $V''$ transversalement
   192	revient à dire que $L_1\cap\ldots \cap L_s$ coupe $\bar{V}{}''$
   193	transversalement. En posant $\bar{V}=$ adhérence de $\bar{V}{}''$, on
   194	est ramené à l'assertion (i) pour $\bar{V}$, qui implique (ii) avec
   195	$V''$ au lieu de $V'$. Il reste donc à considérer le point $P$, si
   196	c'est un point lisse de $V$. Mais il est évident que si
   197	$\{H_1,\ldots,H_s\}$ est assez général et $P$ est lisse sur $V$, alors
   198	$Y$ coupe $V$ transversalement en $P$, d'où le résultat.
   199	
   200	\section{Construction des bons voisinages}\etiquette[3]{XI.sec3}
   201	
   202	\begin{enonce*}{Définition 3.1}\etiquette[3.1]{XI.def3.1}
   203	On appelle fibration élémentaire un morphisme de schémas
   204	${f:X\rightarrow S}$ qui peut être plongé dans un diagramme commutatif
   205	$$
   206	\xymatrix@=2cm{
   207	X  \ar[r]^{j} \ar[rd]_{f}& \bar{X} \ar[d]^{\bar{f}}& Y \ar[l]_{i}  \ar[ld]^{g}\\
   208	               &    S   &      
   209	}
   210	$$
   211	satisfaisant\pageoriginale aux conditions suivantes:
   212	\begin{enumerate}
   213	\renewcommand{\theenumi}{\roman{enumi}}
   214	\renewcommand{\labelenumi}{(\theenumi)} 
   215	% \item $i$ est une immersion fermée.
   216	\item $j$ est une immersion ouverte dense dans chaque fibre, et
   217	  $X=\bar{X}-Y$. 
   218	
   219	\item $\bar{f}$ est lisse et projectif, à fibres géométriques
   220	  irréductibles et de dimension $1$.
   221	
   222	\item $g$ est un revêtement étale, et chaque fibre de $g$ est non-vide.
   223	\end{enumerate}
   224	\end{enonce*}
   225	
   226	
   227	On vérifie facilement qu'un tel plongement de $X$ est unique s'il
   228	existe, à isomorphisme canonique près.
   229	
   230	\begin{enonce*}{Définition 3.2}\etiquette[3.2]{XI.def3.2}
   231	On appelle bon voisinage relatif à $S$ un $S$-schéma $X$ tel qu'il
   232	existe des $S$-schémas
   233	$$
   234	X=X_n,\ldots,X_0=S
   235	$$
   236	et des fibrations élémentaires $f_i:X_i\rightarrow X_{i-1}$, % $f_i$ est un $S$-morphisme
   237	$i=1,\ldots, n$.
   238	\end{enonce*}
   239	
   240	\begin{enonce*}{Proposition 3.3}\etiquette[3.3]{XI.prop3.3}
   241	\footnote{Ce résultat s'étend sans difficulté à une assertion
   242	    semi-locale.}
   243	Soit $k$ un corps algébriquement clos, $X/\Spec k$ un
   244	    schéma lisse, et $x\in X$ un point rationnel. Il existe un ouvert
   245	    de $X$ contenant $x$ qui est un bon voisinage (relatif à $\Spec k$).
   246	\end{enonce*}
   247	
   248	Démonstration. Prenons $X$ irréductible. Par récurrence sur $\dim
   249	X=n$, il suffit de trouver un voisinage $U$ de $x$ dans $X$ et une
   250	fibration élémentaire $f:U\rightarrow V$, avec $V$ lisse et de
   251	dimension $n-1$. En effet, il existera un voisinage $V'$ de $v=f(x)$
   252	qui est un bon voisinage, et on pourra prendre $U'=U\cap \sisi{f^{-1}(V')}{V'}$ comme bon
   253	voisinage de $x$.
   254	
   255	On peut supposer $X\subset \EE^r$ affine. Soit $X_0$
   256	l'adhérence de $X$ dans $\PP^r$. Soit $\bar{X}$ le normalisé de
   257	$X_0$ et $Y=\bar{X}-X$, avec la structure induite
   258	réduite. Soit\pageoriginale de plus $S\subset \bar{X}$ le
   259	sous-ensemble fermé des points singuliers. On a $S\subset Y$ et
   260	\begin{align*}
   261	\dim \bar{X} &= \dim X=n{\quoi},\\
   262	\dim Y &= n-1{\quoi},\\
   263	\dim S &\leq n-2{\quoi}.
   264	\end{align*}
   265	
   266	Plongeons $\bar{X}$ dans un espace projectif $\PP^N$, au moyen
   267	d'un faisceau inversible $L$ qui soit de la forme $M^{\otimes r}$ avec
   268	$M$ très ample et $r\geq 2$. Il existe alors des hyperplans % linéaires
   269	$H_1,\ldots, H_{n-1}$ de $\PP^N$, où $H_i$ est l'ensemble des
   270	zéros de
   271	$$
   272	\sum^N_{\nu =0}a_{i\nu} x_{\nu} =0{\quoi},
   273	$$
   274	qui contiennent $x$ et tels que l'intersection $L=H_1\cap\ldots \cap
   275	H_{n-1}$ soit de dimension $N-n+1$ et coupe $\bar{X}$ et $Y$
   276	transversalement (\hyperref[XI.thm2.1]{2.1}). % voir le commentaire du théorème 2.1
   277	 %Pour l'intersection avec $\bar{X}$: $S\cap L=∅$ pour $L$ assez
   278	 %général, donc $\bar{X}\cap L=(\bar{X} - Y)\cap L$ et
   279	 %$\bar{X}- Y$ est lisse.
   280	 %Pour l'intersection avec $Y$: $Y$ est réduit et $k$ algébriquement clos,
   281	 %donc chaque composante irréductible est génériquement lisse, et $L$ assez
   282	 %général ne rencontre pas le lieu singulier de $Y$.
   283	L'intersection $\bar{X}\cap L$ est une courbe
   284	lisse et connexe (cela résulte du théorème de Bertini), et $Y\cap L$
   285	est de dimension $0$. Choisissons de plus un autre hyperplan
   286	$$
   287	H_0:\sum^N_{\nu=0}a_{0 \nu}x_{\nu}=0
   288	$$
   289	tel que $H_0$ coupe $\sisi{\bar{X}}{X}\cap L$ transversalement et $H_0\cap
   290	Y\cap L=∅$. % Il faut de plus choisir $x\notin H_0$ pour que
   291	                    % l'application $f'$ plus bas soit bien définie en $x$
   292	
   293	Considérons la projection $\PP^N\rightarrow \PP^{n-1}$
   294	obtenue en « introduisant les coordonnées projectives »
   295	$$
   296	y_i=\sum^N_{\nu =0}a_{i_{\nu}}x_{\nu}{\quoi}.
   297	$$
   298	C'est une application rationnelle définie en dehors du centre de
   299	projection $C=H_0\cap\ldots \cap H_{n-1}$. Soit $\epsilon
   300	:P'\rightarrow \PP^N$ l'éclatement de $C$, tel qu'on ait un
   301	diagramme
   302	$$
   303	\xymatrix@=1.5cm{
   304	\PP^{N}  \ar[rd]& &P'\ar[ld]^{\pi}\ar[ll]_{\epsilon}\\
   305	               &\PP^{n-1} &
   306	}
   307	$$
   308	où\pageoriginale $\pi$ est un morphisme. Soit $\bar{X}{}'\subset P'$
   309	l'image inverse « propre » de $\bar{X}$, c'est-à-dire,
   310	l'adhérence de $\epsilon^{-1}(\bar{X}-(\bar{X}\cap C))$. Puisque par
   311	hypothèse $C$ coupe $\bar{X}$ transversalement, le morphisme
   312	$\bar{X}{}'\rightarrow \bar{X}$ est un éclatement de l'ensemble fini
   313	$\bar{X}\cap C$. Soient $X'=X-\bar{X}\cap C$, qui s'identifie aussi à
   314	un sous-schéma ouvert de $\bar{X}{}'$, et $Y'=\bar{X}{}'-X'$ le
   315	sous-schéma fermé avec la structure induite réduite. On a un diagramme
   316	$\DD$ de morphismes
   317	$$
   318	\xymatrix@=2cm{
   319	  X'  \ar[r]^{j} \ar[rd]_{f'}& \bar{X}{}'  \ar[d]^{\bar{f}{}'} & Y'\ar[l]_{i} \ar[ld]^{g'}     &\\
   320	                           & \PP^{n-1} &\qquad ,   
   321	}
   322	$$
   323	et je dis qu'il existe un voisinage $V$ de $v=f'(x)$ tel que la
   324	restriction de $\DD$ à $V$ satisfasse aux conditions (i), (ii),
   325	(iii) de (\hyperref[XI.def3.1]{3.1}), donc que $f'_{\mid V}$ soit une fibration
   326	élémentaire. Cela achèvera la démonstration.
   327	
   328	La condition (i) sera triviale. Pour (ii), notons que $\bar{X}\cap L$
   329	est une courbe lisse par hypothèse, et on voit immédiatement qu'on a
   330	un morphisme biunivoque ${\bar{f}{}'}^{-1}(v)\rightarrow \bar{X}\cap L$ % morphisme biunivoque? 
   331	induit par $\epsilon$. Donc $\left({\bar{f}{}'}^{-1}(v)\right)_{\red}\rightarrow
   332	\bar{X}\cap L$ est bijectif. Pour vérifier que $\bar{f}{}'$ est lisse
   333	% supprimer les \left et \right ?
   334	au-dessus d'un voisinage de $v$, il suffit par le lemme de Hironaka
   335	(\SGA 1 II 2.6 %\ref
   336	\footnote{et \EGA IV 5.12.10.}) de vérifier qu'il
   337	est lisse au point générique de ${\bar{f}{}'}^{-1}(v)$. En ce point,
   338	$\bar{X}{}'$ est isomorphe à $\bar{X}$, et le morphisme est lisse parce
   339	que $L$ coupe $\bar{X}$ transversalement.
   340	
   341	Il reste à démontrer que $g'$ est étale dans un voisinage de $v$
   342	(puisque $Y$ est de dimension $n-1$, il est évident que chaque fibre
   343	de $g'$ est non-vide). On a $Y'=\epsilon^{-1}(Y)\amalg D_1
   344	\amalg\ldots \amalg D_r$, où $\bar{X}\cap C=\{P_1,\ldots,P_r\}$ et
   345	$D_i$ est l'éclatement\pageoriginale de $P_i$ dans $\bar{X}$. Chaque $D_i$
   346	s'identifie à $\epsilon^{-1}(P_i)$, et s'envoie isomorphiquement sur
   347	$\PP^{n-1}$\sisi{. De plus, la projection}{} est définie en chaque point de
   348	$Y$%car $Y\cap C=∅$
   349	, et le morphisme
   350	induit sur $Y$ n'est autre que $g'_{\mid \epsilon^{-1}(Y)}$. Il est étale
   351	au-dessus de $v$, car $L$ coupe $Y$ transversalement, donc $g'$ est
   352	étale au-dessus de~$v$.
   353	
   354	\section{Le théorème de comparaison}\etiquette[4]{XI.sec4}
   355	
   356	\setcounter{subsection}{-1}
   357	\subsection{}\etiquette[4.0]{XI.subsec4.0}
   358	Soit $X$ un schéma localement de type fini sur $\Spec \CC$ et
   359	considérons la topologie usuelle sur l'espace $X(\CC)$ des
   360	points rationnels de $X$. On va noter par $X_{\cl}$ le site des
   361	\emph{isomorphismes locaux} $f:U\rightarrow X(\CC)$,
   362	c'est-à-dire, des morphismes $f$ d'espaces topologiques ayant la
   363	propriété suivante:
   364	
   365	\sisi{pour}{Pour} chaque $x\in U$ il existe un voisinage ouvert $U_x$ tel que
   366	$f_{\mid U_x}$ est un homéomorphisme de $U_x$ sur un voisinage ouvert de
   367	$f(x)$. 
   368	
   369	Une famille $\{U_{\nu}\rightarrow U\}$ de morphismes de $X_{\cl}$ est
   370	dite \emph{couvrante} si $U$ est la réunion des images des $U_{\nu}$.
   371	
   372	Puisqu'une immersion ouverte est un isomorphisme local, on a un
   373	morphisme (inclusion) de sites (IV\quad ) %manque une référence
   374	\footnote{Plus précisément, l'inclusion de catégories
   375	  $\Ouv(X(\CC))\rightarrow X_{\cl}$ définit un morphisme de
   376	  sites (4.1) en sens inverse.}
   377	\begin{equation*}
   378	\delta:X_{\cl}\longrightarrow X(\CC)\qquad
   379	(\ast){\quoi}.\tag{4.1}
   380	\end{equation*}\etiquette[4.1]{XI.eq4.1}
   381	Or pour chaque isomorphisme local $U\rightarrow X(\CC)$ il
   382	existe d'après la définition un « recouvrement » de $U$ par des
   383	ouverts de $X(\CC)$, et on déduit de (\quad) que les topos %manque une référence
   384	associés aux deux sites sont équivalents (par $\delta_{\ast}$). En
   385	particulier on peut remplacer $X(\CC)$ par $X_{\cl}$ pour le
   386	calcul de la cohomologie usuelle.
   387	
   388	Soit maintenant $f:X'\rightarrow X$ un morphisme étale de
   389	schémas. Alors $f(\CC):X'(\CC)\rightarrow X(\CC)$
   390	est un isomorphisme local. Cela résulte immédiatement du\pageoriginale critère
   391	jacobien (\SGA 1 II 4.10) et du théorème des fonctions implicites. Le
   392	foncteur $X'\mapsto X'(\CC)$ donne donc un morphisme de sites
   393	\begin{equation*}
   394	\epsilon:X_{\cl}\longrightarrow X_{\et}{\quoi},\tag{4.2}
   395	\end{equation*}\etiquette[4.2]{XI.eq4.2}
   396	où $X_{\et}$ est le site étale.
   397	
   398	Rappelons les résultats suivants \footnote{Qui figurent dans
   399	  l'exposé cité en note de bas de page, p.~2.}: %identifier les exposés
   400	                                %précédemment pour identifier cet exposé
   401	                                %mentionné dans la note?
   402	
   403	\begin{enonce*}{Théorème 4.3}\etiquette[4.3]{XI.thm4.3}
   404	\begin{enumerate}
   405	\renewcommand{\theenumi}{\roman{enumi}}
   406	\renewcommand{\labelenumi}{(\theenumi)}
   407	\item $X$ est connexe et non-vide si et seulement si $X(\CC)$
   408	  est connexe et non-vide.
   409	
   410	\item Le foncteur $\epsilon$ est pleinement fidèle.
   411	
   412	\item (« Théorème de Grauert-Remmert »). Le foncteur $\epsilon$
   413	  induit une équivalence de la catégorie des revêtements finis de
   414	  $X(\CC)$, muni de sa topologie habituelle d'espace localement
   415	  compact, avec la catégorie des revêtements étales de $X$.
   416	\end{enumerate}
   417	\end{enonce*}
   418	
   419	Nous accepterons (i) comme connu. Il résulte par exemple de (GAGA) %\ref
   420	\cite{4}. L'assertion (ii) résulte facilement de (i). En effet, pour
   421	vérifier l'assertion de surjectivité contenue dans (ii), soit
   422	$\varphi:X'(\CC)\rightarrow X''(\CC)$ un
   423	$X(\CC)$-morphisme. On peut supposer tous les schémas connexes
   424	et affines, donc séparés. Alors le graphe $\Gamma$ de $\varphi$ est
   425	une composante connexe de $\fprod{X'}{X''}{X}(\CC)$, d'où
   426	$\Gamma=Y(\CC)$ pour une certaine composante connexe $Y$ de
   427	$\fprod{X'}{X''}{X}$, par (i). Alors $Y$ est le graphe d'un morphisme
   428	$f:X'\rightarrow X''$ qui induit $\varphi$, \ie la projection
   429	$Y\rightarrow X'$ est un isomorphisme: en effet
   430	$Y(\CC)\rightarrow X(\CC)$ est un isomorphisme, d'où on
   431	tire facilement qu'il en est de même de $Y\rightarrow X'$ (qui est
   432	étale, radiciel, surjectif !). Pour (iii) tout revient d'après (ii) à
   433	démontrer que chaque revêtement fini de $X(\CC)$ est de la
   434	forme $X'(\CC)$ où  $X'\rightarrow X$\pageoriginale est étale.
   435	
   436	Pour cela, on se réduit immédiatement par descente (\SGA 1 IX 4.7) au %\ref
   437	cas $X$ connexe et normal. Le problème est local sur $X$ (pour la
   438	topologie de Zariski), et on peut donc supposer $X$ affine. Soient
   439	$f:X\rightarrow \EE^n$ un morphisme fini, et $T'\rightarrow
   440	X(\CC)$ un revêtement fini connexe. On constate immédiatement
   441	que le morphisme composé $\eta':T'\rightarrow
   442	\EE^n(\CC)$ est un « revêtement analytique » dans
   443	le sens de (\cite{2} \S \hyperref[XI.sec2]{2}, Def.~\hyperref[XI.sec3]{3}). 
   444	D'après (\cite{2} \S \hyperref[XI.sec2]{2}, Satz 8) \sisi{$\eta'$}{} s'étend à un
   445	« revêtement analytique » $\eta:T\rightarrow
   446	\PP^n(\CC)$, et il suffit de démontrer que $\eta$ est
   447	\sisi{induit}{induite} par un morphisme de schémas $f:Y\rightarrow \PP^n$
   448	convenable, avec $Y$ normal. Or le théorème fondamental (\cite{2} \S 13
   449	Satz 42) affirme qu'il y a une structure canonique analytique normale
   450	sur $T$ telle que $\eta$ induise un morphisme d'espaces
   451	analytiques. Puisque alors $\eta$ est un morphisme fini, l'image
   452	directe du faisceau structural sur $T$ est un faisceau $\mathfrak{a}$
   453	d'algèbres analytiques cohérentes sur $\PP^n$. Il résulte de
   454	\sisi{(GAGA)}{GAGA} \cite{4} % \ref
   455	que $\mathfrak{a}$ est le faisceau analytique associé à un
   456	faisceau algébrique $A$ sur $\PP^n$, et on prend
   457	$Y=\sheaf{\Spec}A$.
   458	
   459	Nous pouvons maintenant démontrer le résultat suivant:
   460	
   461	\begin{enonce*}{Théorème 4.4}\etiquette[4.4]{XI.thm4.4}
   462	Soit $X$ un schéma lisse sur $\Spec \CC$.
   463	\begin{enumerate}
   464	\renewcommand{\theenumi}{\roman{enumi}}
   465	\renewcommand{\labelenumi}{(\theenumi)}
   466	\item Il y a une équivalence entre la catégorie des faisceaux de
   467	  torsion localement constants constructibles sur $X_{\et}$ et la
   468	  catégorie des faisceaux de torsion localement constants, à fibres
   469	  finies, sur $X_{\cl}$, l'équivalence étant donnée par les foncteurs
   470	  quasi-inverses ${}_{\ast}$ et ${}^{\ast}$.
   471	
   472	\item Soit $F$ un faisceau de torsion localement constant, à fibres
   473	  finies, sur $X_{\cl}$, et notons aussi par $F$ le faisceau
   474	  $\epsilon_{\ast}F$ induit sur $X_{\et}$. Alors
   475	$$
   476	R^q\epsilon_{\ast}F=0\quad \mbox{ si } q>0{\quoi}.
   477	$$
   478	
   479	\item Avec\pageoriginale les notations de {\rm(i)}, le morphisme
   480	  canonique
   481	$$
   482	H^q(X_{\et},F)\longrightarrow H^q(X_{\cl},F)
   483	$$
   484	est bijectif pour chaque $q$. En particulier,
   485	$$
   486	H^q(X_{\et}, \ZZ/\sisi{n\ZZ}{n})\xrightarrow{\sim} H^q(X_{\cl},\ZZ/\sisi{n\ZZ}{n})
   487	$$
   488	pour chaque $q$.
   489	\end{enumerate}
   490	\end{enonce*}
   491	
   492	Démonstration. Pour (i), notons que les faisceaux en question sont
   493	ceux qui sont représentables par des revêtements étales de $X$ (\resp
   494	par des revêtements finis de $X(\CC)$). Puisqu'il y a une
   495	équivalence entre les catégories de ces revêtements
   496	(\hyperref[XI.thm4.3]{4.3}), l'assertion 
   497	(i) en résulte immédiatement. L'isomorphisme de (iii) est conséquence
   498	de (ii) et de la suite spectrale de Leray V 5.2 % \ref + 5.3 au lieu de 5.2?
   499	pour $\epsilon$. Il
   500	reste donc à démontrer (ii).
   501	
   502	Or $R^q\epsilon_{\ast}F$ est le faisceau associé au préfaisceau
   503	$\Rr^qF,$ où $(\Rr^qF)(X')=H^q(X'_{\cl},F)$ pour
   504	$X'\rightarrow X$ étale. Pour démontrer que $R^q\epsilon_{\ast}F=0$ si
   505	$q>0$, on doit donc démontrer le \sisi{lemme}{} suivant:
   506	
   507	\begin{enonce*}{Lemme 4.5}\etiquette[4.5]{XI.lem4.5}
   508	Soit $\xi\in H^q(X_{\cl},F)$ et $x\in X_{\cl}=X(\CC)$. Il
   509	existe un morphisme étale $X'\rightarrow X$ dont l'image contient $x$
   510	et tel que l'image de $\xi$ dans $H^q(X'_{\cl},F)$ soit nulle.
   511	\end{enonce*}
   512	
   513	Démonstration du lemme. Récurrence sur $n=\dim X$; le problème est
   514	local sur $X$ pour la topologie étale, \sisi{}{et} on peut donc supposer $F$
   515	constant, et par dévissage on peut même supposer que
   516	$F=\ZZ/\sisi{n\ZZ}{n}$. De plus, en remplaçant $X$ par un voisinage
   517	ouvert de Zariski de $x$, on peut (\hyperref[XI.prop3.3]{3.3}) supposer que $X$ admet une
   518	fibration élémentaire (\hyperref[XI.def3.1]{3.1}) $f:X\rightarrow S$. Reprenons les
   519	notations de (\hyperref[XI.def3.1]{3.1}): par\pageoriginale calcul direct, on voit que
   520	$R^q{j}_{\ast}F=0$ si $q>1$, que $R^1j_{\cl\!\ast}F$ est
   521	l'extension par $0$ d'un faisceau constant sur $Y_{\cl}$, et enfin que
   522	$j_{\cl\!\ast}F$ est le faisceau constant $\ZZ/\sisi{n\ZZ}{n}$ sur
   523	$\bar{X}$. De plus, on peut calculer $R^q
   524	\bar{f}_{\cl\!\ast}(j_{\cl\!\ast}F)$ fibre par fibre. De la suite
   525	spectrale de Leray
   526	$$
   527	E^{pq}_2=R^p\bar{f}_{\cl\!\ast}(R^qj_{\cl\!\ast}F)\Longrightarrow
   528	R^{p+q}f_{\cl\!\ast}F
   529	$$
   530	on déduit qu'on peut de même calculer $R^{p+q}f_{\cl\!\ast}F$ fibre
   531	par fibre, et donc que 
   532	
   533	$R^0f_{\cl\!\ast}\ZZ/\sisi{n\ZZ}{n}=\ZZ/\sisi{n\ZZ}{n}$,
   534	
   535	$R^1f_{\cl\!\ast}\ZZ/\sisi{n\ZZ}{n}$ est un faisceau localement constant de
   536	  torsion sur $S_{\cl}$,
   537	
   538	$R^qf_{\cl\!\ast}\ZZ/\sisi{n\ZZ}{n}=0$ si $q>1$.
   539	
   540	La suite spectrale de Leray
   541	
   542	$E^{pq}_2=H^p(S_{\cl}, R^qf_{\cl\!\ast}F)\Longrightarrow
   543	H^{p+q}(X_{\cl},F)$
   544	
   545	\noindent
   546	se réduit ainsi à la suite exacte
   547	\begin{equation*}
   548	\ldots\longrightarrow H^q(S_{\cl},f_{\cl\!\ast}F)\longrightarrow
   549	H^q(X_{\cl},F)\longrightarrow H^{q-1}(S_{\cl},
   550	R^1f_{\cl\!\ast}F)\longrightarrow\tag{$\ast$}
   551	\end{equation*}
   552	Soit $s$ l'image de $x$ dans $S$. Par hypothèse de récurrence et grâce
   553	à (i), il existe pour chaque classe $\eta\in H^q(S_{\cl},G)$ ($G$
   554	localement constant de torsion, à fibres finies) un « voisinage
   555	étale » $S'\rightarrow S$ de $s$ tel que l'image de $\eta$ dans
   556	$H^q(S'_{\cl},G)$ soit nul. Or les foncteurs $R^qf_{\cl\!\ast}$
   557	commutent évidemment aux changements de base étales, et si
   558	$S'\rightarrow S$ est un voisinage étale de $s$, alors
   559	$X'=\fprod{X}{S'}{S}\rightarrow X$ est un voisinage étale de
   560	$x$. Lorsque $q\geq 1$ on termine donc grâce à la suite exacte
   561	$(\ast)$ en prenant d'abord $G=R^1f_{\cl\!\ast}F$ et
   562	ensuite\pageoriginale $G=f_{\cl\!\ast}F$. %La preuve ici est terminée, la
   563	                                %suite est une remarque pour le cas $q=1$
   564	Lorsque $q=1$, on note que
   565	$H^1(X_{\cl},\ZZ/n \ZZ)$ (\resp
   566	$H^1(X,\ZZ/n\ZZ)$) classifie les revêtements étales %plutôt parler de
   567	                                %revêtements finis pour $X_\cl$ ?
   568	principaux de $X_{\cl}$ (\resp $X$), de groupe
   569	$\ZZ/n\ZZ$, donc en vertu de \hyperref[XI.thm4.3]{4.3} (iii) l'application
   570	$H^1(X,\ZZ/n\ZZ)\rightarrow
   571	H^1(X_{\cl},\ZZ/n\ZZ)$ est bijective, d'où aussitôt le
   572	résultat d'effacement \hyperref[XI.lem4.5]{4.5} pour $q=1$.
   573	
   574	\begin{enonce*}[remark]{Variante 4.6}\etiquette[4.6]{XI.var4.6}
   575	On peut aussi démontrer (\hyperref[XI.lem4.5]{4.5}) de la manière élégante suivante, due à
   576	Serre: on remarque qu'un bon voisinage (\hyperref[XI.def3.2]{3.2}) connexe $X$ satisfait aux
   577	deux conditions suivantes:
   578	\begin{enumerate}
   579	\renewcommand{\theenumi}{\roman{enumi}}
   580	\renewcommand{\labelenumi}{(\theenumi)}
   581	\item $\pi_n(X(\CC))=0$ si $n>1$.
   582	
   583	\item $\pi_1(X(\CC))$ est une extension successive de groupes
   584	libres de type fini.
   585	\end{enumerate}
   586	\end{enonce*}
   587	
   588	En effet, si $X\rightarrow S$ est une fibration élémentaire où $S$ est
   589	un bon voisinage, on peut supposer (i) et (ii) vrai{s} pour $S$ par
   590	récurrence. Or $X(\CC)\rightarrow S(\CC)$ est un espace
   591	fibré localement trivial, comme on voit aisément, à fibres isomorphes
   592	à $X_0(\CC)$, où $X_0$ est une fibre arbitraire de
   593	$X/S$. Or $X_0$ est une courbe non-singulière connexe \emph{non
   594	complète} et il est bien connu que cela implique que
   595	$\pi_n(X_0(\CC))=0$ si $n>1$, et que
   596	$\pi_1(X_0(\CC))$ est un groupe libre. % \ref
   597	Alors (i) et (ii)
   598	sont des conséquences de la suite exacte d'homotopie
   599	$$
   600	\longrightarrow \pi_0(X_0(\CC))\longrightarrow
   601	\pi_n (X(\CC))\longrightarrow \pi_n
   602	(S(\CC))\longrightarrow\ldots
   603	$$
   604	
   605	Donc $X(\CC)$ est un espace $K(\pi,1)$, et cela implique que
   606	chaque classe de cohomologie $\xi\in H^n(X(\CC),F)$ devient
   607	nulle sur le revêtement universel de $X(\CC)$. Il suffit de
   608	démontrer que $\xi$ devient nulle déjà sur un revêtement fini $X'$ de
   609	$X$, ce qui revient au même que de démontrer que
   610	$\pi_1=\pi_1(X(\CC))$ est\pageoriginale un \emph{bon groupe}
   611	(\cf CG 1 16), %\ref
   612	 c'est-à-dire que le morphisme de $\pi_1$ dans son
   613	complété $\sisi{\widehat}{}{\pi}_1=\varprojlim_N \pi_1/N$ (où $N$ parcourt l'ensemble des
   614	sous-groupes invariants % 'invariant' = 'distingué'?
   615	de $\pi_1$ d'indice fini) induit une
   616	bijection.
   617	$$
   618	H^q(\widehat{\pi}_1,M)\xrightarrow{\sim} H^q(\pi_1,M)
   619	$$
   620	pour tout $\pi_1$-module fini $M$ continu (où le symbole de gauche est
   621	la cohomologie de $\widehat{\pi}_1$ en tant que groupe profini). Or
   622	une extension successive de groupes libres de type fini est un bon
   623	groupe (\cf CG I p.~15-16, exc.~1,2), d'où le résultat. % \ref
   624	
   625	\begin{thebibliography}{99}
   626	\bibitem{1} GIRAUD J., « Analysis Situs », Séminaire
   627	Bourbaki, Exposé nᵒ 256.
   628	
   629	\bibitem{2} GRAUERT H\sisi{.}{} et REMMERT R., « Komplexe Räume »,
   630	Math, Ann. Bd. 136, p. 245 (1958).
   631	
   632	\bibitem{3} HIRONAKA H., « Resolution of singularities of an
   633	algebraic variety over a field of characteristic zero » Annals of
   634	Math. Vol. 39, nᵒ 1, 2 (1964).
   635	
   636	\bibitem{4} SERRE J.P., « Géométrie Algébrique et Géométrie
   637	Analytique », Annales de l'Institut Fourier, t. VI, p. 1-12 (1956).
   638	\end{thebibliography}
   639	
   640	
   641	\ifx\danslelivre\undefined
   642	\end{document}
   643	\fi
